\documentclass{beamer}
\input{tbook-deck}
\title{Two moves, diamonds and squares}
\subtitle{Bajnok--Beatty Conjectures 4.1 and 1.6}
\author{Haobo Ma \and Wenlin Zhang}
\institute{}
\date{}
\begin{document}
% Manuscript: §1, Theorem 1; §6, Corollary 10
\begin{frame}
\titlepage
\note{This talk explains version 2 of the paper by Haobo Ma and Wenlin Zhang. Credit the original questions and attaining constructions to B. Bajnok and W. K. Beatty. Begin with the two-move game before introducing the algebra.}
\end{frame}
% Manuscript: §1–§6, proof dependency order
\begin{frame}{The question and the route}
\begin{enumerate}
\item Can two moves cover the whole group?
\item Why does a square force congruences?
\item What is the largest possible noncyclic group?
\end{enumerate}
\medskip
Covering $\to$ relation lattice $\to$ rigidity or residue counts.
\note{The two proof routes share the coefficient lattice. Near the upper limit, use rigidity. For arbitrary first invariant factor, use residue counts.}
\end{frame}
\section{Classification and mechanism}
% Manuscript: §1, definition of s-spanning pair
\begin{frame}{An eight-position example}
In $\mathbb Z/2\mathbb Z\times\mathbb Z/4\mathbb Z$, take
\[a=(1,1),\qquad b=(0,1).\]
\begin{center}
\begin{tabular}{c|cccc}
First coordinate & \multicolumn{4}{c}{Second coordinate $0,1,2,3$}\\\hline
$0$ & $0$ & $b$ & $2b$ & $-b$\\
$1$ & $a-b$ & $a$ & $a+b$ & $-a$
\end{tabular}
\end{center}
At most two moves; backwards moves and repetition allowed.
\note{Read the entries as positions in clock arithmetic. Each expression has coefficient cost at most two. In particular $-a=(1,3)$, because the first coordinate is modulo two.}
\end{frame}
% Manuscript: §1, definition; §5, Theorem 1
\begin{frame}{The precise question}
Let $G_k=\mathbb Z/2\mathbb Z\times\mathbb Z/(2k)\mathbb Z$, of order $4k$.
\[\forall g\in G_k\quad \exists x,y\in\mathbb Z:\quad
 g=xa+yb,\quad |x|+|y|\le s.\]
Require $a\ne b$ and $k,s\ge1$.
\begin{block}{A radius-four surprise}
$k=7$: $28$ positions, every pair fails.\\
$k=8$: $32$ positions, a suitable pair succeeds.
\end{block}
\note{The signed coefficients count net uses of the moves. Noncyclic means that no single element generates the group. The next theorem proves the contrast between twenty-eight and thirty-two; it is not a finite search claim.}
\end{frame}
% Manuscript: §1, Theorem 1
\begin{frame}{Conjecture 4.1: the classification}
\begin{bookresult}{thm-classification}
For $k,s\ge1$, an $s$-spanning pair exists exactly when:
\begin{enumerate}
\item $s$ odd, $2k\le s^2-1$;
\item $s$ even, $2k\le s^2-s$;
\item $s\equiv0\pmod4$, $k$ even,\\
$s^2-s<2k\le s^2$;
\item $s\equiv2\pmod4$, $k\equiv2\pmod4$,\\
$s^2-s<2k\le s^2$.
\end{enumerate}
\end{bookresult}
\note{The source is Conjecture 4.1 on page ten of arXiv:2406.04045v1. Preserve the strict lower edge of the upper strip. Radius one permits no positive $k$. Bajnok and Beatty had the size bound and lower regimes; the decisive remaining step is the upper-strip congruence obstruction.}
\end{frame}
% Manuscript: §2, relation lattice and Lemma 2
\begin{frame}{From moves to a diamond and a lattice}
\[D_s=\{(x,y)\in\mathbb Z^2:|x|+|y|\le s\}.\]
\[f(x,y)=xa+yb,\qquad K=\ker f.\]
\begin{itemize}
\item $|D_s|=1+4(1+\cdots+s)=2s^2+2s+1$.
\item Coefficients differing by a vector in $K$ agree.
\item Covering means every coset of $K$ meets $D_s$.
\end{itemize}
Moreover $K\subseteq2\mathbb Z^2$, but $K\not\subseteq4\mathbb Z^2$.
\note{A lattice here is the set of all integer relations that return to the origin. Diamond counting alone is too coarse. Reduction modulo twice the group shows kernel vectors have even coordinates; lifting $(1,0)$ rules out all coordinates being divisible by four.}
\end{frame}
% Manuscript: §2, Lemma 2
\begin{frame}{The diamond becomes a square}
\[T(u,v)=(u+v,u-v),\qquad M=T^{-1}K.\]
\[|x|+|y|=\max\{|x+y|,|x-y|\}.\]
\begin{bookresult}{lem-square}
Every $M$-coset meets $\{0,\ldots,s-1\}^2$.
\[ [\mathbb Z^2:M]=2k\le s^2.\]
\end{bookresult}
Integer parity determines the required translation.
\note{The written lemma covers all positive radii. Even radii require the translate by $(1,0)$ so that the sums and differences are odd. Odd radii give even sums and differences. The transformation has image of index two, halving the index. A rotated picture alone is not the argument.}
\end{frame}
% Manuscript: §3, Lemma 3
\begin{frame}{Almost full means rigid}
\begin{bookresult}{lem-rigidity}
$n\ge2$, $0\le\delta{}<n$, index $n^2-\delta$; every coset meets the $n\times n$ square. Then
\[L=\mathbb Z(n,-\epsilon B)+\mathbb Z(-A,\epsilon n),\]
\[0\le A,B<n,\quad \epsilon=\pm1,\quad AB=\delta.\]
\end{bookresult}
The squeeze: $n(n-1)<d=n^2-\delta$, and $n^2<2d$.
\note{Cover the column just beyond the right boundary to find relation vectors. Their determinant is a positive multiple of the index and at most $n^2$, so equals the index. If neither reaches a full side, it is at most $n(n-1)$, a contradiction. The appendix includes the zero-deficit case.}
\end{frame}
% Manuscript: §3, Lemma 4
\begin{frame}{Where the congruence comes from}
\begin{bookresult}{lem-congruences}
If $s=2t$, $k=2t^2-r$, $0\le r<t$, then
\[2\mid r,\qquad t\text{ odd}\Longrightarrow4\mid r.\]
\end{bookresult}
\begin{itemize}
\item Rigidity gives $AB=2r$.
\item Even coordinate sums force $A=2u,B=2v$: $r=2uv$.
\item If $t,u,v$ are odd, the transformed basis lies in $4\mathbb Z^2$: contradiction.
\end{itemize}
\note{The small deficit makes the rigidity lemma apply. If $t$ is odd and $r$ has just one factor of two, $u,v$ are both odd. Transforming the basis then puts the whole kernel in four times the integer plane, contradicting the earlier kernel fact.}
\end{frame}
% Manuscript: §4, Lemma 5; §5, Lemma 6
\begin{frame}{Constructing the converse}
\begin{bookresult}{lem-family}
For $t\ge1$, $0\le j<t$, with $t$ or $j$ even,
$\mathbb Z/2\mathbb Z\times\mathbb Z/(4(t^2-j))\mathbb Z$
has a $2t$-spanning pair.
\end{bookresult}
\[V=(2(t-j),-2(t+j)),\]
\[W=(2(t-1),2(t+1)).\]
Use $r=2j$ in the upper strip; consecutive-coordinate pairs cover the lower ranges.
\note{The determinant is $8(t^2-j)$ and the entry gcd is two, giving the required quotient by Smith normal form. Reduce transformed coefficients to a least-sum nonnegative representative to prove covering. Handle $t=1,j=0$ separately. The reading gives the lower-range integer formulas in Lemma 6.}
\end{frame}
\section{Maximum noncyclic order}
% Manuscript: §6, Corollary 10
\begin{frame}{Conjecture 1.6: all noncyclic groups}
\begin{bookresult}{cor-global}
For $s\ge3$, the maximum order is:
\[N(s)=2s^2\quad (2s+1\text{ prime});\]
\[N(s)=2s^2+2s-\frac{p^2-1}{2}\quad (2s+1\text{ composite}),\]
where $p$ is the smallest prime divisor of $2s+1$.
\end{bookresult}
$s=3$: $N(3)=18$;\quad $s=4$: $N(4)=36$.
\medskip
Attainment: Bajnok and Beatty. This paper completes the upper bound for arbitrary pairs.
\note{A maximum theorem needs both an upper bound and attainment. The corresponding formal theorem includes both. At radius four, fixing the first factor at two gives thirty-two positions, while allowing first factor three gives thirty-six.}
\end{frame}
% Manuscript: §6, Lemma 9
\begin{frame}{The sparsest residue class limits size}
Write $G=\mathbb Z/c\mathbb Z\times\mathbb Z/(ck)\mathbb Z$.
Each coefficient residue class modulo $c$ must cover $k$ elements.
\medskip
For odd $c$, sum and difference are invertible modulo $c$; the count is
\[E_iE_j+O_iO_j.\]
At $s=4,c=3$:
\begin{itemize}
\item Residue $2$: $-4,-1,2$, giving $(E,O)=(2,1)$;
\item Residue $0$: $-3,0,3$, giving $(1,2)$.
\end{itemize}
Thus $k\le2\cdot1+1\cdot2=4$, and order $\le3^2\cdot4=36$.
\note{Each coefficient residue class corresponds to one coset of $cG$, with $k$ elements. Since the kernel lies in $c\mathbb Z^2$, different coefficient residues cannot represent one element. Sum and difference must have equal parity, giving even times even plus odd times odd.}
\end{frame}
% Manuscript: §6, proof of Corollary 10
\begin{frame}{Why the smallest prime factor?}
Exceeding $2s^2$ requires
\[s=ct+\frac{c-1}{2},\qquad c\mid2s+1.\]
Then
\[|G|\le2s^2+2s-\frac{c^2-1}{2}.\]
The penalty grows with $c$: the smallest prime factor $p$ gives the best value.
\medskip
If $2s+1$ is prime, $c=2s+1$ forces $t=0,k=0$, which is impossible.
\note{All other remainder cases are bounded by the baseline. In the composite case the least prime divisor gives the largest candidate upper bound, and the existing Bajnok–Beatty construction with $t=((2s+1)/p-1)/2$ attains it.}
\end{frame}
% Manuscript: §6, Corollary 11 and perfect spanning sets paragraph
\begin{frame}{Equality and the boundary of the result}
\begin{bookresult}{cor-universal}
For $s\ge3$, $|G|\le2s^2+2s-4$. Equality holds exactly for
\[s\equiv1\pmod3,\]
\[m=(2s^2+2s-4)/3,\]
\[G\simeq\mathbb Z/3\mathbb Z\times\mathbb Z/m\mathbb Z.\]
\end{bookresult}
Small radii: $N(2)=9$; no noncyclic example at radius one.
\medskip
Perfect pairs force cyclicity. The corrected full multiplier-shape assertion is not proved here.
\note{This equality-type classification retains written-proof status. A formal numerical upper bound alone would not cover this whole statement. Perfect representation forces order $2s^2+2s+1$, above the noncyclic bound. Cyclicity is distinct from a full classification of the pair shape.}
\end{frame}
\appendix
\section{Supporting details}
% Manuscript: §3, proof of Lemma 3
\begin{frame}{Appendix: the rigidity squeeze}
Put $d=n^2-\delta>n(n-1)$. Covering the outside column gives
\[(p,-b),(r,c)\in L,\quad1\le p,r\le n.\]
Without a short horizontal vector, minimal $b,c<n$ satisfy $b+c\le n$.
\[0<pc+rb\le n(b+c)\le n^2<2d.\]
The determinant is a multiple of $d$, hence equals $d$.
\[b+c=n,\qquad p=n\ \text{or}\ r=n.\]
A short horizontal vector instead forces $\delta=0$.
\note{The integer intervals $[0,n-1-b]$ and $[c,n-1]$ must cover every height. If $b+c<n$, or both horizontal components are below $n$, the determinant is at most $n(n-1)$, contradicting the index. Taking the difference of the vectors yields the asserted basis.}
\end{frame}
% Manuscript: §6, Lemma 9
\begin{frame}{Appendix: the three-part residue bound}
\begin{bookresult}{lem-odd}
For odd $c\ge3$, $s=ct+r$, $0\le r<c$,
\[k\le\begin{cases}
2t^2,&0\le r<(c-1)/2,\\
2t(t+1),&(c-1)/2\le r<c-1,\\
2t(t+1)+1,&r=c-1.
\end{cases}\]
For arbitrary pairs, including $t=0$.
\end{bookresult}
\note{Choose residue $s+1$, whose representatives alternate in parity. In the first range there are $2t$. In the middle range residues $s+1,s+2$ have $2t+1$ representatives with opposite majority parity. In the last range, use the same residue twice to obtain $t^2+(t+1)^2$.}
\end{frame}
% Manuscript: Bibliography; §6, attaining constructions
\begin{frame}{Sources and attribution}
\small
Haobo Ma and Wenlin Zhang, paper version 2.\\
\href{https://doi.org/10.5281/zenodo.23931731}{Version DOI: 10.5281/zenodo.23931731}.
\medskip

B. Bajnok and W. K. Beatty: original questions, earlier ranges, attaining constructions.\\
\href{https://arxiv.org/abs/2406.04045v1}{arXiv:2406.04045v1}; journal DOI: 10.2140/cnt.2025.14.1.
\medskip

R. Dougherty and V. Faber: quotient-lattice viewpoint, Proposition 2.\\
DOI: 10.1137/S0895480100372899.
\medskip

Full proofs and focused contribution requests are in the reading editions.
\note{Use the original Bajnok–Beatty preprint for conjecture numbering. Distinguish original mathematical authors, prior constructions and formalization contributors. Exact statement correspondence determines the evidence status; a paper-level source link does not verify every statement.}
\end{frame}
% Manuscript: statement evidence for cited results
\FormalVerification
\end{document}
